\chapter{Linear programming}
\label{chap:pl}

\section{What a linear model represents}
Linear programming (LP) allocates scarce resources among activities when
both the objective and the constraints are linear
\cite{dantzig1963linear,bertsimas1997introduction}. A variable may represent
units produced, hours assigned, or quantities transported. Its meaning
and unit must be established before writing coefficients.

\begin{definicion}
A continuous linear model has the form
\[
 \operatorname{opt}_{x\in\R^n} c^\top x,\qquad
 a_i^\top x\ \{\leq,=,\geq\}\ b_i\quad(i=1,\ldots,m),\qquad
 \ell_j\leq x_j\leq u_j.
\]
The operator denotes minimization or maximization. A bound may be absent;
a free variable admits both positive and negative values.
\end{definicion}

Linearity assumes proportionality and additivity: doubling an activity
doubles its resource use and contribution, and contributions add without
interaction effects. A continuous formulation allows divisibility. The data
are also deterministic over the modeled horizon. Discounts, start-ups,
indivisible batches, or variable yields require a different formulation. Validating
a file checks its structure and mathematical rules; it does not establish that
these assumptions accurately represent the organization.

\section{Example: product mix}
A plant manufactures two divisible products, \(X_1\) and \(X_2\). Their unit
contributions are 50 and 60 monetary units. Each unit uses two resources as
shown in the following table; capacities cover the same period as production.
\begin{table}[htbp]
\centering
\caption{Product mix data}
\begin{tabular}{lrrr}
\toprule
Resource & Per unit of \(X_1\) & Per unit of \(X_2\) & Capacity \\
\midrule
Resource 1 & 2 & 3 & 180 \\
Resource 2 & 3 & 2 & 150 \\
\bottomrule
\end{tabular}
\end{table}
The model is
\[
\begin{aligned}
 \max\quad &50X_1+60X_2\\
 \text{subject to}\quad &2X_1+3X_2\leq180,\\
 &3X_1+2X_2\leq150,\\
 &X_1,X_2\geq0.
\end{aligned}
\]
Capacities are hard constraints: exceeding them is not allowed. Minimum
demand requirements and fixed costs are not included; attributing them to this
result would change the problem.

\subsection{Reproducible input and execution}
The complete model is in \path{manual/examples/lp_sample.json}.
\par\noindent The keys
\texttt{sense}, \texttt{variableNames}, \texttt{objectiveCoefficients}, and
\texttt{constraints} define the objective and rows. Each row includes a name,
coefficients in variable order, a relation, and \texttt{rhs}. Bounds are
\texttt{lowerBounds} and \texttt{upperBounds}; \texttt{null} in an upper bound
means that there is no explicit maximum. The example uses
\texttt{continuous} types. Run from the repository root:
\begin{lstlisting}
swift run qsb solve-json manual/examples/lp_sample.json --backend native
\end{lstlisting}
This specific command returns \texttt{objectiveValue} and the
\texttt{variableValues} dictionary. It does not include a sensitivity report.

On macOS, open the normalized JSON in the
\texttt{Linear / Integer Programming} family. The \texttt{Definition} view
shows \texttt{Model}, \texttt{Variables and objective}, and the constraints.
To enter it manually, select \texttt{Maximize}, create two variables with
coefficients 50 and 60, type \texttt{Continuous}, a lower bound of zero, and no
upper bound; then enter the two rows with relation \(\leq\). Use
\texttt{Validate} and correct the diagnostics before pressing \texttt{Run LP}.
The editor shortcuts are Command--Shift--V to validate and Command--R to
solve. \texttt{Copy Table} and \texttt{Paste Table} transfer tables to and from
a spreadsheet. Preserve the column order and headers;
review the model after pasting. In narrow windows, the editor may require
horizontal scrolling.

\subsection{Result and interpretation}
\begin{table}[htbp]
\centering
\caption{Linear optimum verified with QSB}
\begin{tabular}{lr}
\toprule
Quantity & Value \\
\midrule
\(X_1\) & 18 \\
\(X_2\) & 48 \\
Total contribution & 3780 \\
\bottomrule
\end{tabular}
\end{table}
The solution view shows the objective, variables, and constraint details.
Check feasibility numerically:
\[
 2(18)+3(48)=180,\qquad3(18)+2(48)=150.
\]
Both resources are fully used; their slacks \(b_i-a_i^\top x\) are zero.
An active constraint holds as an equality at the optimum; a positive slack
indicates unused capacity. This interpretation depends on the relation:
for \(a_i^\top x\geq b_i\), the surplus is \(a_i^\top x-b_i\).

The feasible region is a polygon with vertices \((0,0)\), \((50,0)\),
\((18,48)\), and \((0,60)\). Their objective values are 0, 2500, 3780, and 3600.
This independent geometric check confirms the result for the two-variable
problem. In higher dimensions, enumerating vertices by hand is impractical.

\section{Duality and sensitivity: what can be inferred}
For this maximization problem with constraints \(Ax\leq b\), \(x\geq0\), the dual is
\[
 \min\ 180y_1+150y_2,\qquad
 2y_1+3y_2\geq50,\quad3y_1+2y_2\geq60,\quad y_1,y_2\geq0.
\]
Since both production variables are positive at the optimum, both dual
constraints hold as equalities. They yield
\(y_1=16\), \(y_2=6\), and \(180(16)+150(6)=3780\). The dual values
certify an upper bound equal to the achieved objective.

As long as this basis remains optimal, one additional unit of resource 1
is worth 16 monetary units, and one of resource 2 is worth 6. These are local
marginal values, not prices valid for every expansion. For example, changing
only the first capacity to \(180+\delta\) moves the intersection to
\[
 X_1=18-\tfrac25\delta,\qquad X_2=48+\tfrac35\delta.
\]
Both quantities remain nonnegative for
\(-80\leq\delta\leq45\), the interval over which the objective is
\(3780+16\delta\). At the endpoints the basis structure changes;
outside the interval the problem must be solved again. This derivation is an
academic calculation for the example: QSB currently displays neither shadow
prices nor sensitivity ranges in the specific LP output or its native view.

\section{Method, backends, and limitations}
The \texttt{native} backend uses two-phase simplex: it first seeks a feasible
solution and then improves the objective. It is deterministic and educational,
using floating-point arithmetic. A mathematical optimum does not imply
symbolic precision. Widely differing coefficient scales can impair
calculations; use consistent units and check residuals.

\texttt{validate} runs the checks without computing an optimum.
\texttt{external} requests HiGHS, installed independently and discovered
through \texttt{QSB\_HIGHS\_PATH} or \texttt{PATH}; if it is unavailable, the
request fails explicitly. There is no silent fallback to the native backend.
The educational engine does not provide the preprocessing tools or scale
guarantees of an industrial engine. To exchange the model with another
optimizer, you can generate MPS:
\begin{lstlisting}
swift run qsb export-mps manual/examples/lp_sample.json > /tmp/qsb-lp.mps
\end{lstlisting}
MPS preserves names through a deterministic adaptation and expresses bounds
and domains; JSON remains QSB's normalized representation.

Distinguish three situations: an invalid model contains inconsistent data;
an infeasible model has no decision satisfying all constraints; an unbounded
model admits unlimited improvements in the objective. Adding an arbitrary
bound to make an error disappear changes the model's meaning. Review the
units, signs, and limits that the case actually imposes.

\section{Exercises}
\begin{enumerate}
 \item Increase the capacity of resource 1 to 190. Before running, predict
 \(X_1=14\), \(X_2=54\), and objective 3940 using the sensitivity analysis above;
 confirm the constraints with QSB.
 \item Reduce the capacity of resource 2 to 120 and solve. Explain why a
 marginal valuation should not be extrapolated indefinitely.
 \item Add \(X_1\geq60\). Identify the contradiction with resource 2 and
 compare a structurally valid model with an infeasible problem.
 \item Change the variables to integers and study the next chapter.
 Why does the optimum remain unchanged in the original example?
\end{enumerate}
